\documentclass[11pt,a4paper]{article}
\usepackage[top=0.75in, bottom=0.75in, left=1.0in, right=1.0in]{geometry}
\usepackage{graphicx}
\usepackage[sort&compress, super, numbers]{natbib}
\usepackage[export]{adjustbox}
\usepackage[ngerman]{babel}
\usepackage[utf8]{inputenc}
\usepackage[T1]{fontenc}
\renewcommand{\bibname}{References}%names reference list 
% novelty statement, a cover letter, and a title page including an authorship statement, a data statement, the number of words in abstract/main text and the number of figure/table/boxes.
% recommended and opposed reviewers, and recommended and opposed editorial board members with whom they are in conflict. 
\begin{document}

\pagenumbering{gobble}
\noindent \includegraphics[width=0.4\textwidth, right]{letterhead/ubc-logo-2018-fullsig-blue-cmyk.png}
\noindent Dear Dr. Thrall,
\vspace{1.5ex}\\
\noindent Please consider our paper, ``Traits predict forest phenological responses to photoperiod more than temperature'' for publication as a full paper in \emph{Ecology Letters}. 
\vspace{1.5ex}\\ 
\noindent 
Climate change is altering the environmental cues of growth globally, producing shifts in the timing of life history events (phenology), altering resource use and producing novel species interactions. 
%emwJul12 -- oy, the below sentence is really clunky and long! Can you simplify? I would try to slim the whole below paragraph. Maybe just say -- predicting these shifts is critical but has proven difficult because phenology is rarely included in studies of growth strategies (CITES). Yet phenology is foundational to the theory that underlies these spectra \citep{Grime1977}, which clearly predicts early to late phenology should align with acq to cons growth stratgies and relating traits. This prediction has rarely been tested because phenology is often highly variable when measured in natural conditions, experiments and new models may overcome this challenge.  
But to date, we have observed considerable variation in communities' responses to environmental change, likely reflecting differences in species assemblages and the effects of species' other traits on the timing and patterns of growth. In plants, decades of research has found consistent relationships between traits and growth strategies under different environments---giving rise to frameworks like the leaf or wood economic spectra \citep{Wright2004, Chave2009}. Within these frameworks, species fall along a gradient of fast and acquisitive to slow and conservative  strategies. Phenology was foundational to the theory that underlies these spectra \citep{Grime1977}, but has rarely been included because it is often challenging to measure and highly variable across species when measured in natural conditions. 
 \vspace{1.5ex}\\ 
 Here, we paired two massive controlled environment experiments with trait data for over 1400 individual plants across 47 species and eight sites, testing how phenology fits within major frameworks for plant growth. Based on decades of research into the relationship between traits and plant growth strategies, we explored whether early versus late leafout is correlated with acquisitive versus conservative growth strategies, and how  these relationships change between forests across North America. 
 \vspace{1.5ex}\\ 
 %emwJul12 --I feel like I have variously read 5, 4 and 6 traits (abstract, discussion and here) --- make sure it is consistent ... 
We found that phenology fits within major frameworks for plant growth, with early to late budburst correlating with size and leaf metrics. Of the six traits we tested,  four traits related to leafout through photoperiod cues. Despite previous studies finding temperature to be the dominant controller of leafout, we found no support that traits relationships with temperature cues influence leafout timing. These results emerge in addition to trait variation across sites, with most traits varying across our latitudinal gradients and/or differences between our eastern and western transects of North America. % , likely reflecting differences in winter intensities across the sites.
\vspace{1.5ex}\\ 
Our results provide the most compelling evidence of trait-phenology relationships with implications for how woody plant species will shift with climate change. Photoperiod is generally considered the weakest cue of leafout, but its shared importance with other species traits and shaping growth strategies suggest it may pose a major constraint on how species adapt to warmer climates than previous thought. Accounting for these phenology-trait relationships will improve our ability to accurately forecast future changes in species growth strategies and species interactions without forest communities. %emwJul12 -- without? I might re-write this last sentence as: Accounting for these phenology-trait relationships can improve forecasts of forest ecosystem and community change with warming, but have larger implications for how we understand phenology and growth strategies more broadly.  
\vspace{1.5ex}\\ 
\noindent All authors contributed to this work and approve this version for submission. The manuscript is 4014 words with a 149 word summary, and 4 figures and is not under consideration elsewhere. We hope you find it suitable for publication in \emph{Ecology Letters}, and look forward to your response. 
%\noindent We recommend the following reviewers: Dr. Meredith Zettlemoyer, Dr. Mason Heberling, Dr. Jason Fridley, Dr. Rong Yu, and Dr. Ameila Caffara. 
\vspace{1.5ex}\\
\noindent Sincerely, \\
\includegraphics[scale=.4]{letterhead/shot.png} \\ 
\noindent Deirdre Loughnan\\
\noindent Sentinels of Change Postdoctoral Fellow\\ 
\noindent Hakai Institute $|$ Department of Zoology\\
\noindent University of British Columbia
\newpage
\vspace{-5ex}

\begingroup
\renewcommand{\section}[2]{}
\bibliographystyle{refs/bibstyles/nature.bst}% 
\bibliography{refs/biblio.bib}
\endgroup
\newpage


\end{document}

%\begingroup
%\renewcommand{\section}[2]{}%
%\renewcommand{\chapter}[2]{}% for other classes
%\bibliography{mybibfile}
%\endgroup